\documentclass[onecolumn, draftclsnofoot,10pt, compsoc]{IEEEtran}
\usepackage{graphicx}
\usepackage{url}
\usepackage{setspace}

\usepackage{geometry}
\geometry{textheight=9.5in, textwidth=7in}

\usepackage{pgfgantt}
\definecolor{grey}{rgb}{0.57, 0.64, 0.69}

% 1. Fill in these details
\def \CapstoneTeamName{The Cleverly Named Team}
\def \CapstoneTeamNumber{20}
\def \GroupMemberOne{Cade Raichart}
\def \GroupMemberTwo{Daniel Kato}
\def \GroupMemberThree{Nathan Shepherd}
\def \CapstoneProjectName{ISS Barometer App }
\def \CapstoneSponsorCompany{NASA}
\def \CapstoneSponsorPerson{Don Pettit}

% 2. Uncomment the appropriate line below so that the document type works
\def \DocType{	%Problem Statement
				% Requirements Document
				%Technology Review
				Design Document
				%Progress Report
				}

\newcommand{\NameSigPair}[1]{\par
\makebox[2.75in][r]{#1} \hfil 	\makebox[3.25in]{\makebox[2.25in]{\hrulefill} \hfill		\makebox[.75in]{\hrulefill}}
\par\vspace{-12pt} \textit{\tiny\noindent
\makebox[2.75in]{} \hfil		\makebox[3.25in]{\makebox[2.25in][r]{Signature} \hfill	\makebox[.75in][r]{Date}}}}
% 3. If the document is not to be signed, uncomment the RENEWcommand below
%\renewcommand{\NameSigPair}[1]{#1}

%%%%%%%%%%%%%%%%%%%%%%%%%%%%%%%%%%%%%%%
\begin{document}
\begin{titlepage}
    \pagenumbering{gobble}
    \begin{singlespace}
    	\includegraphics[height=4cm]{coe_v_spot1}
        \hfill
        % 4. If you have a logo, use this includegraphics command to put it on the coversheet.
        \includegraphics[height=4cm]{NASA}
        \par\vspace{.2in}
        \centering
        \scshape{
            \huge CS Capstone \DocType \par
            {\large\today}\par
            \vspace{.5in}
            \textbf{\Huge\CapstoneProjectName}\par
            \vfill
            {\large Prepared for}\par
            \Huge \CapstoneSponsorCompany\par
            \vspace{5pt}
            {\Large\NameSigPair{\CapstoneSponsorPerson}\par}
            {\large Prepared by }\par
            Group\CapstoneTeamNumber\par
            % 5. comment out the line below this one if you do not wish to name your team
            %\CapstoneTeamName\par
            \vspace{5pt}
            {\Large
                \NameSigPair{\GroupMemberOne}\par
                \NameSigPair{\GroupMemberTwo}\par
                \NameSigPair{\GroupMemberThree}\par
            }
            \vspace{20pt}
        }
        \begin{abstract}
        % 6. Fill in your abstract
This Requirements Document is meant as a repeatable and detailed description of using the iPad Air 2 Barometer to create an app for displaying barometer readings.
This app will be used to measure atmospheric pressure in emergency situations aboard the International Space Station.
This document will also serve to be an agreement that both our team and our client will make, and it will be used to evaluate our product at the end of Spring term.
Within this document is the description, constraints and dependencies that will be used for this app.
The Specific requirements that must be completed are laid out within the user, hardware and software interfaces sections.
The System features entails the more in depth features of the pressure readings, recording button, graph and the settings pages.
        \end{abstract}
    \end{singlespace}
\end{titlepage}
\newpage
\pagenumbering{arabic}
\tableofcontents
% 7. uncomment this (if applicable). Consider adding a page break.
%\listoffigures
%\listoftables
\clearpage

\section{Introduction}
The \CapstoneProjectName will provide crew members of the International Space Station with a readout of barometric pressure, as read from the iPads internal sensor.
This application will supplement the current method of air pressure reading, the MANOVACUMETER, a mechanical barometer.
As our application is targeted toward the crew members, it needs to be highly readable, clear, and good at quickly conveying data.
Along with ease of use, the need for readability and efficiency of use are the main decision factors in our production, as well as our design.
Listed below are the final product requirements of this product; version one will be provided at the end of December 2017, version two will be provided before the end of May.

\begin{itemize}
\item[V1:] Data Screen with initial pressure, current pressure, start/stop button, time stamp, pressure change (measured in sec / mmHg)
\item[V1:] Graph Screen with plot pressure as a function of time
\item[V1:] Additional settings menu that includes plot scale, scroll functionality, orientation, decimal point, as well as other other units of measure for pressure
\item[V1:] All items will be optimized for efficient reading and clear interface
\item[V2:] Graph will have pinch to zoom functionality, scaling options (to be set in settings)
\item[V2:] Graph will be able to be exported as a csv, excel readable file
\item[V2:] Graph can be compared to t-Res table, to measure time remaining
\item[Optional:] Research can be done to identify low pressure accuracy and resistance
\end{itemize}

As approved by our client, these are the goals of our project, with unlisted items being out of scope.
This may be changed in the future, depending on available time, for instance, the research portion may not be feasible, due to time restrictions.
Our version one application will be given to the client, who will then be able to give it to the ISS crew members.
Although they are the intended target user, giving them version one will allow our team to improve upon any issues that the find in using the application.
This means that our version two application may have changes not mentioned in the list, as they may be specific to the users and will only be apparent after use.
This design document will detail our choices from the individual Technology Reviews, and will provide details related to our decision making, with regards to the mentioned requirements.
\cite{udefaultApple}
\section{Glossary of Terms}
\begin{itemize}
\item[--] csv: A file format readable by Excel and other spread sheet software
\item[--] iOS: The internal operating system of Apple products, specifically iPads
\item[--] iPad: A touch screen tablet device with barometer capabilities
\item[--] ISS: International Space Station
\item[--] MANOVACUMETER: The current mechanical air pressure gauge aboard the ISS
\item[--] mmHg: Millimeters of Mercury, the standard pressure unit
\item[--] NASA: National Aeronautics and Space Administration
\item[--] t-Res: The chart that NASA uses to calculate residual time left on station
\end{itemize}

\section{Data Storage}

\section{Language}

\section{UI Library}

\section{Graph Interface}

\section{Barometer Implementation}

\section{Organization of UI}

\section{Chart Library}

\section{Development Framework}

\section{Color Palette}

\section{Conclusion}

\newpage
\bibliographystyle{IEEEtran}
\bibliography{bibliography}
\end{document}
